% year: תשפ"ב
% type: מבחן
% semester: א
% moed: א
% number: 1
% points: 17
% categories: continuity, derivatives
% source: exams/מבחנים/תשפב/מועד א תשפב.pdf

\begin{question}
נגדיר פונקציה על ידי
\[ f(x)=\begin{cases} x^3 \sin(\frac{5}{x}) & x\neq 0,\\ 0 & x=0 \end{cases} \]
\begin{enumerate}
  \item (6 נק') האם $f$ רציפה בנקודה $x=0$ ? נמקו.
  \item (6 נק') האם $f$ גזירה בנקודה $x=0$ ? נמקו.
  \item (5 נק') האם הנגזרת $f'(x)$ רציפה בנקודה $x=0$ ? נמקו.
\end{enumerate}
\end{question}

\begin{hint}
"אפסה כפול חסומה": $\left|\sin(\frac{5}{x})\right|\le 1$ לכל $x\neq 0$.
\end{hint}

\begin{hint}
בסעיף ב' חשבו את הנגזרת ב-$0$ לפי ההגדרה (גבול מנת ההפרשים). בסעיף ג' חשבו את $f'(x)$ עבור $x\neq0$ לפי כללי הגזירה ובדקו את הגבול כאשר $x\to 0$.
\end{hint}

\begin{solution}
\begin{enumerate}
  \item \textbf{כן.} לכל $x\neq 0$ מתקיים
  \[ 0\le |f(x)| = |x|^3\left|\sin\left(\tfrac{5}{x}\right)\right| \le |x|^3 \xrightarrow[x\to0]{} 0, \]
  ולכן לפי משפט הסנדוויץ' $\lim_{x\to0} f(x)=0=f(0)$. כלומר $f$ רציפה ב-$0$.
  \item \textbf{כן.} לפי הגדרת הנגזרת:
  \[ f'(0)=\lim_{h\to0}\frac{f(h)-f(0)}{h}=\lim_{h\to0}\frac{h^3\sin(5/h)}{h}=\lim_{h\to0} h^2\sin\left(\tfrac{5}{h}\right). \]
  כאן $h^2\to0$ ו-$\sin(5/h)$ חסומה, ולכן (אפסה כפול חסומה) הגבול שווה $0$. לכן $f$ גזירה ב-$0$ ו-$f'(0)=0$.
  \item \textbf{כן.} עבור $x\neq0$, לפי כלל המכפלה וכלל השרשרת:
  \[ f'(x)=3x^2\sin\left(\tfrac{5}{x}\right)+x^3\cos\left(\tfrac{5}{x}\right)\cdot\left(-\tfrac{5}{x^2}\right)=3x^2\sin\left(\tfrac{5}{x}\right)-5x\cos\left(\tfrac{5}{x}\right). \]
  שני המחוברים הם מהצורה "אפסה כפול חסומה" ($3x^2\to0$, $5x\to0$ ו-$|\sin|,|\cos|\le1$), ולכן
  \[ \lim_{x\to0}f'(x)=0=f'(0). \]
  מכאן ש-$f'$ רציפה בנקודה $x=0$.
\end{enumerate}
\end{solution}
